% Drafted for this port (no source docs/SECTION_7_DRAFT.md exists).
% Synthesizes only claims already established in Sections 1, 3-6; no new
% claims are introduced here, per the task's own instruction.
\section{Conclusion}
\label{sec:conclusion}

We set out to answer a question the graph-ERC literature's own experimental contract had left unexamined: does conversation-graph machinery retain measurable value once the utterance encoder beneath it is itself fine-tuned? Under a residual attribution construction in which every architectural component --- multimodal fusion, speaker-state recurrence, relational pair-state memory, a shift-aware auxiliary objective, temporal attention --- provably initializes as an identity function on a frozen foundation (Section~\ref{sec:residual-construction}), the answer is a clean and uniform \emph{no}. On frozen features the same instrument reproduces the literature's premise exactly, with a positive, multi-seed-consistent gain of $+0.0356 \pm 0.0106$ (Section~\ref{sec:meld-frozen}); on a fine-tuned foundation, at the locked recipe and even at a context width of zero, the identical machinery's paired contribution is null-to-negative (Section~\ref{sec:meld-finetuned}--\ref{sec:k-sweep}). The gain does not partially survive fine-tuning; it disappears at the first point fine-tuning is introduced, before any conversational context is even added to the encoder's input.

A pre-registered trial gave the strongest single component of that machinery, relational memory, its best conceivable case to overturn this verdict: a corpus of long dyadic dialogues, a protocol and label chosen in advance specifically because they favor relationship modeling, on a foundation retrained with the same locked recipe (Section~\ref{sec:iemocap}). The result is not merely another null; it is a null in the diagnostic direction, more negative on IEMOCAP than on MELD, the opposite of the pre-registered prediction. Yet the mechanism is demonstrably not broken: relational memory measurably improves an auxiliary shift-detection task on both corpora even as it degrades the primary emotion classification (Section~\ref{sec:dissociation}). The component learns real conversational-dynamics structure; a fine-tuned foundation has simply already captured the portion of that structure the emotion label needs, leaving the graph pathway with signal that is genuine but redundant, or actively distracting, for the task at hand.

These findings are not an indictment of the graph-ERC literature's original results, which were real and reproducible on the foundations available when they were reported (Section~\ref{sec:related-work}). They are an indictment of treating a frozen-feature comparison as sufficient evidence of an architectural component's value once fine-tuning a strong encoder costs minutes on a single GPU rather than the specialized infrastructure it once required. Our counsel to the field is accordingly a change of default, not a moratorium on conversational structure: attribute new architectural components against a competitively fine-tuned encoder baseline, with an instrument that separates a component's true contribution from the incidental cost of its own re-training, before treating a frozen-feature gain as evidence the component itself --- rather than the representational gap beneath it --- is what is doing the work. We release the residual attribution methodology, the full experimental code, and the complete pre-registered decision history, including this paper's own negative results, so that this standard can be applied, checked, and extended by others.
